% Complete finite-chain splice proof of Kamp III.2, Lemma 6.

\section{Kamp's fixed-initial-point lemma}
\label{sec:annex-p09-kamp-iii2-lemma6}

\begin{annexitem}{annex:p9.kamp.iii2.lemma6}{thm:p09-kamp-iii2-lemma6}{Both initial-point cases and every joining-interval position}
The source's proof begins on printed p.~79 and ends immediately before
Lemma~7 on p.~80. We use the fixed occurrence word and its proved
realization semantics, and supply the finite switch-index argument
and the full interval verification
\cite[ch.~III, lemma~6, pp.~79--80]{Kamp1968}.

\begin{proof}
Fix \(\mathcal M=\langle T,<,(P_i^{\mathcal M})_i\rangle\), a
nonempty strict linear structure, and the same
\(\phi\in\operatorname{Bp}_n(x,y)\), \(n\geq1\), throughout.
Write its word as \((A_0,B_1,A_1,\ldots,B_n,A_n)\).
Let \(a<b<c<d\) and assume \(\phi(a,b)\) and \(\phi(a,d)\).
All point-argument notation has the semantic meaning fixed in the body.
Applying a predicate letter to a point below asserts membership in
that letter's interpretation in \(\mathcal M\).

\paragraph{The forward direction and the exact past witness.}
If \(\phi(a,c)\) holds, choose \(a\) as the initial-point witness
in \(\phi^{\exists<}(c)\). The strict inequality \(a<c\) follows
from the given four-point order. The free-variable replacement in its
definition evaluates \(\phi[v/x]\) at \(v=a,y=c\) exactly as
\(\phi\) at \(x=a,y=c\); bound variables were made fresh.
Hence the guarded existential is true. This direction requires no
restriction on its witness beyond the displayed strict past guard.

For the reverse direction, suppose \(\phi^{\exists<}(c)\) holds.
It supplies some point \(v_0<c\) with \(\phi(v_0,c)\).
The basic-pattern realization equivalence in
\cref{annex:p9.kamp.iii2.lemma1} gives a sequence
\[
                           v_0<\cdots<v_n=c
\]
realizing the whole fixed word. Every position \(v_i\) satisfies
\(A_i\), and each \((v_{i-1},v_i)\) satisfies \(B_i\)
throughout. This includes the final point condition \(A_n(c)\).
There is no assumption comparing \(v_0\) with \(a\); the following
two cases cover all its possible positions.

\paragraph{The first switch index exists within the source range.}
Suppose \(a\leq v_0\). Use \(\phi(a,d)\) to choose a realization
\[
                           a=w_0<\cdots<w_n=d.
\]
Thus \(w_0\leq v_0\), whereas \(v_n=c<d=w_n\).
Let
\[
                    J=\{i\in\{0,\ldots,n\}:w_i\leq v_i\}.
\]
The finite set \(J\) contains 0 and does not contain \(n\).
It has a greatest member \(j\) with \(0\leq j\leq n-1\).
Maximality gives \(w_{j+1}>v_{j+1}\); the strict order in
\(\mathbf v\) then gives
\[
                     w_j\leq v_j<v_{j+1}<w_{j+1}.
\]
In particular \(j\) satisfies the source switch condition
\(w_j\leq v_j<w_{j+1}\). Therefore the set of eligible indices
in \(\{0,\ldots,n-1\}\) is nonempty and has a least member.
Only a maximum of a finite set of integers was used here.
No supremum or infimum in the order \(T\) was assumed.

The construction may choose that least eligible \(k\), or any other
eligible index. In the remainder of this case assume exactly
\[
                      w_k\leq v_k<w_{k+1},
                      \qquad 0\leq k\leq n-1.
\]
We do not assume the additional inequality
\(v_{k+1}<w_{k+1}\) that was available for the existence witness
\(j\). Consequently the following verification proves the actual
source condition for every allowed choice.

\paragraph{The first splice: strict order, all point labels and the old gaps.}
Define the output sequence \(\mathbf r=(r_0,\ldots,r_n)\) by
\[
 r_i=
 \begin{cases}
   w_i,&0\leq i\leq k,\\
   v_i,&k+1\leq i\leq n.
 \end{cases}
\]
Its initial portion is strictly increasing as part of \(\mathbf w\);
its final portion is strictly increasing as part of \(\mathbf v\).
At the join,
\(w_k\leq v_k<v_{k+1}\), so \(r_k<r_{k+1}\).
Thus all \(n+1\) points are in strict order, with
\(r_0=w_0=a\) and \(r_n=v_n=c\).

For every \(i\leq k\), the point \(r_i=w_i\) satisfies \(A_i\);
for every \(i\geq k+1\), the point \(r_i=v_i\) satisfies that
same position's \(A_i\). In particular the initial label comes from
the realization beginning at \(a\), and the final label at \(c\)
comes from the guarded past realization.

For interval indices \(1\leq i\leq k\), the output gap is the
old \((w_{i-1},w_i)\), which satisfies \(B_i\).
For \(k+2\leq i\leq n\), it is the old
\((v_{i-1},v_i)\), which also satisfies its \(B_i\).
The only remaining gap is \((w_k,v_{k+1})\), at index \(k+1\).

\paragraph{The first joining interval: the two old gaps cover every point.}
Take any \(t\) with \(w_k<t<v_{k+1}\).
If \(t<w_{k+1}\), then
\(w_k<t<w_{k+1}\), so the old \(\mathbf w\) gap supplies
\(B_{k+1}(t)\). If \(w_{k+1}\leq t\), use
\(v_k<w_{k+1}\leq t<v_{k+1}\); then the old \(\mathbf v\)
gap supplies \(B_{k+1}(t)\).
Linearity makes these two cases exhaustive, including
\(t=w_{k+1}\). This proves
\[
 (w_k,v_{k+1})\subseteq
                  (w_k,w_{k+1})\cup(v_k,v_{k+1}),
\]
and both sets on the right satisfy the same fixed interval predicate
\(B_{k+1}\).

There is no hidden requirement that an old point label imply an
interval label. If \(w_k<v_k\), the discarded point \(v_k\)
lies inside \((w_k,w_{k+1})\), so that old gap supplies its
newly required \(B_{k+1}\) truth. If
\(w_{k+1}<v_{k+1}\), the discarded point \(w_{k+1}\)
lies inside \((v_k,v_{k+1})\), so the other gap supplies its
\(B_{k+1}\) truth. If \(w_k=v_k\), that point is the new
left endpoint and the open-gap condition does not test it.
If \(w_{k+1}=v_{k+1}\), that point is the new right endpoint
and is likewise excluded. These cover the equalities as well as both
possible relative orders of the two right endpoints.
Every output gap is now verified, so \(\mathbf r\) realizes
\(\sigma(\phi)\) from \(a\) to \(c\). Hence \(\phi(a,c)\).

\paragraph{The second switch index and the other source premise.}
Suppose instead that \(v_0\leq a\). The case needed after the
first branch is \(v_0<a\), but proving the weak case also includes
the equality permitted in the source. Use \(\phi(a,b)\) to choose
\[
                           a=w_0<\cdots<w_n=b.
\]
We have \(v_0\leq w_0\), while \(w_n=b<c=v_n\).
Put
\[
                    K=\{i\in\{0,\ldots,n\}:v_i\leq w_i\}.
\]
Again 0 belongs to the finite set and \(n\) does not. Its greatest
member \(j<n\) satisfies \(v_{j+1}>w_{j+1}>w_j\).
Thus
\(v_j\leq w_j<v_{j+1}\), proving the existence of an eligible
index \(k\in\{0,\ldots,n-1\}\) for this second source condition.
The least eligible index used by the method is therefore defined.
This is another finite integer comparison argument.

\paragraph{The second splice: every point and interval.}
For any eligible \(k\), return the same indexed sequence:
\(r_i=w_i\) for \(i\leq k\), and \(r_i=v_i\) for
\(i\geq k+1\). Its two portions are strictly increasing, and the
switch inequality \(w_k<v_{k+1}\) supplies the strict order
at the join. Its endpoints are \(a,c\), since the retained
prefix contains index 0 and the retained suffix contains index \(n\).

Every point before or at the switch retains its \(A_i\) from
\(\mathbf w\), and every point after it retains its \(A_i\)
from \(\mathbf v\). Thus \(A_0(a)\), \(A_n(c)\) and every
interior point label hold. The gaps at indices \(1,\ldots,k\)
retain the corresponding \(\mathbf w\) intervals and their \(B_i\)
labels. Those at indices \(k+2,\ldots,n\) retain the
corresponding \(\mathbf v\) intervals and their \(B_i\) labels.
For the single new gap, \(v_k\leq w_k\) gives
\[
                         (w_k,v_{k+1})\subseteq(v_k,v_{k+1}).
\]
The right-hand gap of the old \(\mathbf v\) realization supplies
\(B_{k+1}\) at each new gap point. Equality \(v_k=w_k\)
only makes the two left endpoints equal; it creates no closed-endpoint
test. Therefore all realization conditions hold for \(\mathbf r\),
and the basic-pattern equivalence yields \(\phi(a,c)\).

The first case uses the realization ending at \(d\), and the
second uses the one ending at \(b\). Since every point \(v_0\)
of a linear order satisfies \(a\leq v_0\) or \(v_0\leq a\),
these cases prove the reverse implication completely.

\paragraph{Switch boundaries, repeated letters and empty intervals.}
For \(k=0\), the output retains only \(w_0=a\) from its prefix.
There are no old prefix gaps to check; the joining gap has index 1,
and the suffix starts at \(v_1\). For \(k=n-1\), it retains
only \(v_n=c\) from its suffix. There are no old suffix gaps;
the joining gap has index \(n\). Both are within the source's
one-based range \(1,\ldots,n\). The output always has \(n+1\)
points and the original \(n\) intervals.

At \(n=1\), the sole possible switch is \(k=0\). Both
constructions return \((a,c)\); the gap arguments above apply
to its one interval. There is no stage \(n=0\) in the source
basic-pattern grammar. If \(v_0=a\), index 0 is eligible in
both branches: \(w_0=v_0<w_1\) in the first and
\(v_0=w_0<v_1\) in the second. The least-index method returns
the original sequence \(\mathbf v\), already beginning at \(a\).

All interval assertions concern open intervals. If a tested interval
is empty, its universal predicate condition is vacuous. The interval
inclusions prove the required new condition even when one or both
covering gaps are empty. If the relevant \(B_{k+1}\) is false
everywhere, those inclusions force the new gap itself to be empty;
they do not demand a nonexistent point. Repeated predicate letters
cause no change: the same interpretation is used whenever that letter
occurs, and the proof checks each occurrence at its original position.
No truth set or predicate interpretation is altered by the splice.
Adjacent points in an input realization have an empty intervening
gap and impose no interval-predicate test there. If \(a,b\) themselves
are adjacent and \(n>1\), the first realization premise has no witness.
For \(n=1\), that premise permits their empty gap and the proof above applies.

Four strictly ordered points are an explicit hypothesis. If the
domain lacks such a quadruple, that hypothesis has no instance.
If there are not enough points to realize a chosen pattern between
specified endpoints, the corresponding realization premise is false.
The construction never supplies points to repair either absence.
It permits global first and last points whenever the displayed four
points and realizations exist.

\paragraph{Verification of the overlapping-gap example.}
In the first example, \((0,2,3)\) satisfies \(A_0,A_1,A_2\)
at its respective points; its first gap contains only 1, which
belongs to \(B_1\), and its last gap is empty.
The sequence \((0,2,8)\) has the same first two point labels,
has its last point in \(A_2\), and satisfies its last \(B_2\)
gap because \(B_2=T\). Thus both required fixed-start realizations
hold at the chosen four points \(0<3<6<8\).

For \((1,4,6)\), the three points belong respectively to
\(A_0,A_1,A_2\). Its first gap consists of 2 and 3, both in
\(B_1\), and its last gap consists of 5, in \(B_2\).
Hence the guarded past closure at 6 holds. The first source case
uses \(\mathbf w=(0,2,8)\); index 0 satisfies
\(w_0=0\leq v_0=1<w_1=2\), so it is the least eligible index.
The resulting \((0,4,6)\) has the required point labels.
Its first gap is \(\{1,2,3\}\subseteq B_1\), and its second
is \(\{5\}\subseteq B_2\). The old gaps supplying the new
first gap are exactly \(\{1\}\) and \(\{2,3\}\).
In particular \(w_1=2<v_1=4\); this example verifies a
source-eligible switch for which the new gap is not contained in
the first old gap alone.

\paragraph{Verification of the final-label example and the source scope.}
For the one-interval word in the second example, the initial point 0
satisfies \(A_0\), the points 1 and 3 satisfy \(A_1\), and
their open intervals from 0 satisfy \(B_1=T\). Thus the two
premises hold. The point 2 does not satisfy \(A_1\), so
\(\phi(0,2)\) is false. Every realization ending at 2 would
also require \(A_1(2)\), so its guarded past closure is false.
This proves both assertions in that example and shows why the
past-closure condition cannot be replaced merely by an inequality.

The proof uses exactly the source's fixed pattern, two fixed-start
realizations, four strictly ordered points and guarded past
condition. The switch inequalities are unchanged; only the point
indices are shifted from \(1,\ldots,n+1\) to \(0,\ldots,n\).
The finite scan for the least eligible index terminates after at
most the \(n\) indicated candidates. The output uses only the
input witness points, with the same point and interval count.
Every order argument used linearity or finite indexing. Therefore
the result has the source's explicit \(\mathsf{Lin}\) scope.
It uses none of the later interval reductions and no additional
completeness assumption.
\end{proof}
\end{annexitem}
